% Part: first-order-logic 
% Chapter: axiomatic-deduction 
% Section: provability

% verification of properties of provability needed for maximally 
% consistent sets in the completeness chapter.

\documentclass[../../../include/open-logic-section]{subfiles}

\begin{document}

\iftag{FOL}
      {\olfileid[pt]{fol}{axd}{prv}}
      {\olfileid[pt]{pl}{axd}{prv}}

\olsection{Propriedades da \usetoken{S}{derivability}}

\begin{prop}[Monotonicidade] Se $\Gamma \subseteq \Delta$ e $\Gamma \Proves !A$,
então também $\Delta \Proves !A$. \end{prop}

\begin{proof} Todo $\Gamma_0 \subseteq \Gamma$ finito também é um subconjunto
finito de~$\Delta$; portanto, !!a{derivation} de $\Gamma_0 \Proves !A$ também mostra
$\Delta \Proves !A$. \end{proof}

\begin{prop}\ollabel{prop:provability} $\Gamma \Proves !A$ se, e somente se,
$\Gamma\cup \{\lnot!A \}$ é inconsistente. \end{prop}

\begin{prop} \begin{enumerate}
\item \ollabel{prop:provability-contr} Se $\Gamma \Proves !A$ e $\Gamma
\cup \{ !A\} \Proves \lfalse$, então $\Gamma$ é inconsistente.

\item \ollabel{prop:provability-lnot} Se $\Gamma \cup \{!A\}
\Proves \lfalse$, então $\Gamma \Proves \lnot !A$.

\item \ollabel{prop:provability-exhaustive} Se $\Gamma \cup \{!A\}
\Proves \lfalse$ e $\Gamma \cup \{\lnot !A\} \Proves
\lfalse$, então $\Gamma \Proves \lfalse$.

\tagitem{prvOr}{\ollabel{prop:provability-lor-left} Se $\Gamma \cup \{!A\}
\Proves \lfalse$ e $\Gamma \cup \{!B\} \Proves
\lfalse$, então $\Gamma \cup \{!A \lor !B\} \Proves \lfalse$.}{}

\tagitem{prvOr}{\ollabel{prop:provability-lor-right} Se $\Gamma
\Proves !A$ ou $\Gamma \Proves !B$, então $\Gamma
\Proves !A \lor !B$.}{}

\tagitem{prvAnd}{\ollabel{prop:provability-land-left} Se $\Gamma
\Proves !A \land !B$, então $\Gamma \Proves !A$ e
$\Gamma \Proves !B$.}{}

\tagitem{prvAnd}{\ollabel{prop:provability-land-right} Se $\Gamma
\Proves !A$ e $\Gamma \Proves !B$, então $\Gamma
\Proves !A \land !B$.}{}

\tagitem{prvIf}{\ollabel{prop:provability-mp} Se $\Gamma \Proves
!A$ e $\Gamma \Proves !A \lif !B$, então $\Gamma
\Proves !B$.}{}

\tagitem{prvIf}{\ollabel{prop:provability-lif} Se $\Gamma \Proves
\lnot !A$ ou $\Gamma \Proves !B$, então $\Gamma \Proves
!A \lif !B$.}{} \end{enumerate} \end{prop}

\begin{proof} 
\begin{enumerate}
\item Suponha que $\Gamma_0 \cup \{!A\} \Proves \lfalse$.

\begin{derivation} 
1. & $\Gamma_0 \cup \{!A\} \Proves \lfalse$ & hip. \\
2. & $\Gamma_1 \Proves !A$ & hip.\\ 
3. & $\Gamma_0 \cup \Gamma_1 \Proves \lfalse$ & corte \\ 
\end{derivation}

Como $\Gamma_0 \subseteq \Gamma$ e $\Gamma_1 \subseteq \Gamma$,
$\Gamma_0 \cup \Gamma_1 \subseteq \Gamma$; logo, $\Gamma \Proves \lfalse$.

\item Suponhamos que $\Gamma_0 \cup\{\lnot !A\} \Proves \lfalse$.

\begin{derivation} 
1. & $\Gamma_0 \cup\{\lnot !A\} \Proves \lfalse$ & hip. \\ 
2. & $\Gamma_0 \Proves \lnot !A \lif \lfalse$ & Teorema da Dedução\\
3. & $\Gamma_0 \Proves (\lnot !A \lif \lfalse) \lif !A$ & Ax0\\ 
4. & $\Gamma_0 \Proves !A$ & MP 2, 3 \\ 
\end{derivation}

\item Suponhamos que $\Gamma \cup \{ !A \} \Proves \lfalse$ e que
$\Gamma_1 \cup \{\lnot !A \} \Proves \lfalse$.

\begin{derivation} 
1. & $\Gamma_0 \cup\{!A\} \Proves \lfalse$ & hip. \\
2. & $\Gamma_0 \Proves !A \lif \lfalse$ & Teorema da Dedução \\ 
3. & $\Gamma_1 \cup \{\lnot !A \} \Proves \lfalse$ & hip.\\ 
4. & $\Gamma_0 \Proves (!A \lif \lfalse) \lif \lnot !A$ & Ax0\\ 
5. & $\Gamma_0 \Proves \lnot !A$ & MP 2, 4\\ 
6. & $\Gamma_0 \cup \Gamma_1 \Proves \lfalse$ & corte 3, 5 \\
\end{derivation}

Como $\Gamma_0 \subseteq \Gamma$ e $\Gamma_1 \subseteq \Gamma$,
$\Gamma_0 \cup \Gamma_1 \subseteq \Gamma$. Portanto, $\Gamma \Proves \lfalse$.

% prop:provability-lor-left 
\tagitem{defOr}{}{
\iftag{probOr}{Exercício.}{Exercício.}}

% prop:provability-lor-right 
\tagitem{defOr}{}{ \iftag{probOr}{Exercício.}{
Suponhamos que $\Gamma_0 \Proves !A$.

\begin{derivation} 
1. & $\Gamma_0 \Proves !A$ & hip. \\
2. & $\Gamma_0 \Proves !A \lif (!A \lor !B)$ & Ax0\\ 
3. & $\Gamma_0 \Proves !A \lor !B$ & MP 1, 2 \\
\end{derivation}

Portanto, $\Gamma \Proves !A \lor !B$. A demonstração para o caso em que $\Gamma \Proves
!B$ é semelhante.}}

% prop:provability-land-left 
\tagitem{defAnd}{}{ \iftag{probAnd}{Exercício}{
Suponha que $\Gamma_0 \Proves !A \land !B$

\begin{derivation} 
1. & $\Gamma_0 \Proves !A \land !B$ & hip. \\
2. & $\Gamma_0 \Proves (!A\land !B) \lif !A $ & Ax0\\ 
3. & $\Gamma_0 \Proves !A \lor !B$ & MP 1, 2 \\
\end{derivation}

Logo, $\Gamma \Proves !A$. !!{derivation} semelhante mostra que
$\Gamma \Proves !B$.}}

% prop:provability-land-right
\tagitem{defAnd}{}{\iftag{probAnd}{Exercício.}{Exercício.}}

% prop:provability-mp 
\tagitem{defIf}{}{ \iftag{probIf}{Exercício.}{ Suponhamos
tanto que $\Gamma_0 \Proves !A$ quanto que $\Gamma_0 \Proves !A \lif !B $.

\begin{derivation} 
1. & $\Gamma_0 \Proves !A$ & hip. \\
2. & $\Gamma_0 \Proves !A \lif !B $ & hip.\\ 
3. & $\Gamma_0 \Proves !B$ & MP 1, 2\\
\end{derivation}

Como $\Gamma_0 \cup \Gamma_1 \subseteq \Gamma$, isso significa que $\Gamma
\Proves !B$.}}

% prop:provability-lif 
\tagitem{defIf}{}{\iftag{probIf}{Exercício.}{Primeiro,
suponha que $\Gamma \Proves \lnot !A$.
    
\begin{derivation} 
1. & $\Gamma_0 \Proves \lnot !A$ & hip. \\
2. & $\Gamma_0 \Proves \lnot !A \lif (!A \lif !B) $ & Ax0\\ 
3. & $\Gamma_0 \Proves !A \lif !B$ & MP 1,
2\\ \end{derivation}

A demonstração para $\Gamma \Proves !B$ é semelhante:

\begin{derivation} 
1. & $\Gamma_0 \Proves !B$ & hip. \\
2. & $\Gamma_0 \Proves !B \lif (!A \lif !B) $ & Ax0\\ 
3. & $\Gamma_0 \Proves !A \lif !B$ & MP 1, 2\\
\end{derivation}}}

\end{enumerate} 
\end{proof}

\begin{probtag}{probOr,probAnd,probIf} Complete a demonstração de
\olref[fol][axd][prv]{prop:provability}. 
\end{probtag}

\begin{thm}[Generalização] \ollabel{thm:Generalization} Se $\Gamma \Proves
!A$ e $x$ não é livre em nenhuma !!{formula} de $\Gamma$, então $\Gamma
\Proves \lforall[x][!A]$. 
\end{thm}

\begin{proof} Por indução sobre $\mathsf{Thm}_\Gamma$: se $!A$ é um axioma,
$\lforall[x][!A]$ também é. Se $!A\in \Gamma$, então $x$ não é livre em $!A$,
de modo que $!A \lif \lforall[x][!A]$ é um axioma (\textbf{Ax3}), e $\Gamma
\Proves \lforall[x][!A]$ por MP. Agora, suponha que $!A$ decorra por MP porque
$\Gamma \Proves !B$ e $\Gamma \Proves !B \lif !A$. Pela hipótese de
indução, $\Gamma \Proves \lforall[x][!B]$ e $\Gamma \Proves
\lforall[x][( !B \lif !A)]$. Mas, por \textbf{Ax2}, também \[ \Gamma \Proves
\lforall[x][( !B \lif !A)] \lif (\lforall[x][ !B] \lif \lforall[x][!A]), \]
e duas aplicações de MP fornecem $\Gamma \Proves \lforall[x][!A]$, como
desejado. \end{proof}

\begin{thm}\ollabel{thm:WeakGen} (\emph{Generalização fraca sobre constantes})
Se $\Gamma \Proves !A$ e $c$ é !!a{constant} que não ocorre em
$\Gamma$, então existe uma variável $x$ que não ocorre em $!A$ tal que $\Gamma
\Proves \lforall[x][\Subst{!A}{x}{c}]$, e a demonstração não envolve $c$.
\end{thm}

\begin{proof} Seja $!A_1,\ldots,!A_n$ uma demonstração de $!A$ a partir de $\Gamma$, de modo
que $!A=!A_n$. Escolha uma variável $x$ que não ocorra em $!A_1,\ldots,!A_n$ e considere
a nova sequência: $\Subst{!A_1}{x}{c},\ldots,\Subst{!A_n}{x}{c}.$ Essa
sequência é uma demonstração de $\Subst{!A}{x}{c}$ a partir de $\Gamma$. De fato, para cada
$i=1,\ldots,n$: \begin{itemize} \item se $!A$ é um axioma, então
$\Subst{!A_i}{x}{c}$ também é; \item se $!A_i \in \Gamma$, como $c$ não ocorre em
$\Gamma$, temos $\Subst{!A_i}{x}{c} = !A_i \in \Gamma$; \item se $!A_i$
é obtida de $!A_j$ e $!A_j \lif !A_i$, então $\Subst{!A_i}{x}{c}$
decorre por MP de $\Subst{!A_j}{x}{c}$ e $\Subst{(!A_j \lif !A_i)}{x}{c}
= \Subst{!A_j}{x}{c} \lif \Subst{!A_i}{x}{c}$. \end{itemize} É claro
que a !!{constant} $c$ já não ocorre na nova sequência. Seja agora
$\Gamma'$ o conjunto das !!{formula}s de $\Gamma$ que aparecem na
!!{derivation} de $\Subst{!A}{x}{c}$; então $x$ não é livre em nenhuma
!!{formula} de $\Gamma'$ e $\Gamma' \Proves \Subst{!A}{x}{c}$. Por
generalização (\olref{thm:Generalization}), $\Gamma'
\Proves \lforall[x][\Subst{!A}{x}{c}]$; por monotonicidade, portanto, $\Gamma \Proves
\lforall[x][\Subst{!A}{x}{c}]$, como desejado. 
\end{proof}

\begin{explain}
Nosso objetivo é substituir a exigência, em \olref{thm:WeakGen}, de que $x$
não ocorra em $!A$ (de modo algum) pelas exigências mais fracas de que $x$
não seja livre em $!A$ e seja !!{free for} $c$ em $!A$. Isso pode ser
obtido por uma mudança de !!{variable} ligada --- e é isso que
provaremos primeiro.
\end{explain}

\begin{lem}
\ollabel{lem:free-for} 
Se $x$ é livre para $c$ em $!A$ e $y$ não é livre em $!A$, então $x$
é !!{free for} $y$ em $\Subst{!A}{y}{c}$.
\end{lem}

\begin{proof} Se $x$ não é !!{free for} $y$ em $\Subst{!A}{y}{x}$, então
alguma ocorrência livre de $y$ em $\Subst{!A}{y}{c}$ está no escopo de
um quantificador $\lforall[x]$. Mas, por hipótese, $y$ não é livre em $!A$; portanto,
todas essas ocorrências vêm da substituição $\Subst{}{y}{c}$. Assim, alguma
ocorrência de $c$ está no escopo de $\lforall[x]$, e $x$ não é
!!{free for} $c$ em $!A$. 
\end{proof}

\begin{lem}[Mudança de variável ligada]
\ollabel{lem:ChangeBdVar} 
Se $x$ e $y$ não são livres em $!A$ e ambos são !!{free for} $c$
em $!A$, então $\Proves \lforall[x][\Subst{!A}{x}{c}] \equiv
\lforall[y][\Subst{!A}{y}{c}]$ (e a demonstração não envolve $c$).
\end{lem}

\begin{proof} 
Pelas hipóteses, $y$ é !!{free for} $c$ em $!A$ e $x$ não é
livre em $!A$; portanto, pela \olref{lem:free-for} anterior, $y$ é
!!{free for} $x$ em $\Subst{!A}{x}{c}$. Segue que
$\lforall[x][\Subst{!A}{x}{c} \lif \Subst{!A}{x}{c} \Subst{}{y}{x}]$
é um axioma (\textbf{Ax1}). Como $x$ não é livre em $!A$, temos
$\Subst{!A}{x}{c} \Subst{}{y}{x} = \Subst{!A}{y}{c}$, de modo que 
\[
\Proves \lforall[x][\Subst{!A}{x}{c}] \lif \Subst{!A}{y}{c}, 
\] 
e, pelo Teorema da Dedução, $\lforall[x][\Subst{!A}{x}{c}] \Proves
\Subst{!A}{y}{c}$. Como $y$ não é livre em $!A$, também não é livre em
$\lforall[x][\Subst{!A}{x}{c}]$; assim, por Generalização,
$\lforall[x][\Subst{!A}{x}{c}] \Proves \lforall[y][\Subst{!A}{y}{c}]$,
e o Teorema da Dedução fornece novamente $\Proves
\lforall[x][\Subst{!A}{x}{c}] \lif \lforall[y][\Subst{!A}{y}{c}]$.
A demonstração de $\Proves \lforall[y][\Subst{!A}{y}{c}] \lif \lforall[x]
[\Subst{!A}{x}{c}]$ é perfeitamente simétrica; portanto, a conclusão
decorre por Taut.
\end{proof}

\begin{thm}[Generalização forte sobre constantes]
\ollabel{thm:StrongGen} 
Se $\Gamma \Proves !A$, a !!{constant} $c$ não ocorre em
$\Gamma$, e $x$ não é livre em $!A$, mas é !!{free for} $c$ em $!A$,
então $\Gamma \Proves \lforall[x][\Subst{!A}{x}{c}]$. 
\end{thm}

\begin{proof} 
Como $\Gamma \Proves !A$ e $c$ não ocorre em $\Gamma$, pela
Generalização Fraca existe !!a{variable} $y$ que não ocorre em $!A$ tal que
$\Gamma \Proves \lforall[y][\Subst{!A}{y}{c}]$. Como $y$ não ocorre em $!A$
(de modo algum), não é livre em $!A$ e também é !!{free for} $c$ em $!A$;
se, além disso, (por hipótese) $x$ não é livre em $!A$ e é !!{free for} $c$ em
$!A$, então estão satisfeitas as condições para uma mudança de !!{variable} ligada; logo,
$\Proves \lforall[x][\Subst{!A}{x}{c}] \equiv
\lforall[y][\Subst{!A}{y}{c}]$; daí $\Gamma \Proves
\lforall[x][\Subst{!A}{x}{c}]$. 
\end{proof}

\begin{thm} 
\ollabel{thm:provability-quantifiers}
\begin{tagenumerate}{prvEx,prvAll} 
\tagitem{prvEx}{Se $\Gamma \Proves !A(t)$, então $\Gamma \Proves
  \lexists[x][!A(x)]$.}{} 
\tagitem{prvAll}{Se $\Gamma \Proves \lforall[x][!A(x)]$, então $\Gamma
  \Proves !A(t)$.}{}
\end{tagenumerate} 
\end{thm}

\begin{proof} 
\begin{tagenumerate}{prvEx,prvAll} 
%%Need axioms for exists
\tagitem{prvEx}{Exercício.}{}
\tagitem{prvAll} {Suponha que $\Gamma_0 \Proves \lforall[x][!A(x)]$.

\begin{derivation} 
1. & $\Gamma_0 \Proves \lforall[x][!A(x)]$ & hip. \\
2. & $\Gamma_0 \Proves \lforall[x][!A(x)] \lif \Subst{!A}{t}{x}$ & Ax1 \\ 
3. & $\Gamma_0 \Proves \Subst{!A}{t}{x}$ & MP 1, 2 \\ 
\end{derivation}}{}

\end{tagenumerate} 
\end{proof}

\end{document}
